\chapter{A Horizon-Penetrating Integrator}
\label{ch:xks}

The thesis integrates Carter's separated equations in Boyer--Lindquist (BL)
coordinates, in Mino time, in double precision. For metrology that is ideal:
the constants of motion are exact, and the integrator is tested to
$10^{-12}$. For a real-time renderer the same coordinates bring two problems.
The polar axis $\theta = 0, \pi$ is a coordinate singularity. Rays passing
near it needed clamps, a snap of the angular momentum and a step-squeeze
heuristic, and they still left a faint seam down the image's central column.
And the horizon is a coordinate singularity too ($\Delta \to 0$), where
backward-traced captured rays stall. This chapter describes the replacement
used by every live image of the extension: Hamiltonian geodesics in
Cartesian Kerr--Schild coordinates \citep{KerrSchild2009, Visser2007}.

\section{The Kerr--Schild form}

In Kerr--Schild form the Kerr metric is a flat metric plus a null-vector
term \citep{Kerr1963, KerrSchild2009},
\begin{equation}
g_{\mu\nu} = \eta_{\mu\nu} + f\,l_\mu l_\nu, \qquad
f = \frac{2Mr^3}{r^4 + a^2z^2},
\label{eq:ksmetric}
\end{equation}
\begin{equation}
l_\mu = \left(1,\; s\,\frac{rx + s\,ay}{r^2 + a^2},\;
s\,\frac{ry - s\,ax}{r^2 + a^2},\; s\,\frac{z}{r}\right),
\end{equation}
where the BL radius $r$ is the positive root of
$r^4 - (x^2 + y^2 + z^2 - a^2)\,r^2 - a^2z^2 = 0$ and $s = \pm1$ selects the
chart. Because $l$ is null for both $\eta$ and $g$, the inverse needs no
matrix inversion: $g^{\mu\nu} = \eta^{\mu\nu} - f\,l^\mu l^\nu$. The relation
to BL coordinates is
\begin{equation}
\begin{gathered}
x + iy = (r + i s a)\sin\theta\,e^{i\phi_{\rm KS}}, \qquad z = r\cos\theta, \\
t_{\rm KS} = t_{\rm BL} + s\,T(r), \qquad \phi_{\rm KS} = \phi_{\rm BL} + s\,\Phi(r),
\end{gathered}
\end{equation}
with $T = \int 2Mr/\Delta\,dr$ and $\Phi = \int a/\Delta\,dr$. The two charts
differ in which horizon they cover regularly. The \emph{outgoing} chart
($s = -1$) is regular across the past horizon, which is where camera rays
traced backward in time end when they fall in. The \emph{ingoing} chart
($s = +1$) is regular across the future horizon, which an in-falling observer
crosses (Chapter~\ref{ch:xinfall}).

\section{Hamiltonian equations with an analytic gradient}

With $H = \tfrac12 g^{\mu\nu}p_\mu p_\nu$ and $l\cdot p \equiv l^\mu p_\mu$, the
flat part of the inverse metric is constant, so Hamilton's equations reduce
to
\begin{equation}
\frac{dx^\mu}{d\lambda} = \eta^{\mu\nu}p_\nu - f\,l^\mu\,(l\cdot p),
\qquad
\frac{dp_i}{d\lambda} = \tfrac12\,\partial_i\!\left[f\,(l\cdot p)^2\right],
\qquad
\frac{dp_t}{d\lambda} = 0 .
\label{eq:kshamilton}
\end{equation}
The gradient follows in closed form from $\partial_i r = r(r^2x_i +
a^2z\,\delta_{iz})/(r^4 + a^2z^2)$ and the explicit $f$ and $l$. Unlike the
finite-difference gradients common in real-time renderers, this introduces
no step-size-dependent truncation error. Stationarity and axisymmetry make
$E = -p_t$ and $L_z = x\,p_y - y\,p_x$ exact constants. The Carter constant is
recovered by transforming the covector back to BL components, and $H = 0$
measures the null condition.

The double-precision twin (\texttt{physics/src/kerr\_schild.cpp}) integrates
\eqref{eq:kshamilton} with the thesis' adaptive Dormand--Prince scheme.
Cameras build their rays in BL components in the observer's frame
(Chapter~\ref{ch:xobs}), map them to Kerr--Schild components with the
Jacobian of the coordinate change, and lower them with \eqref{eq:ksmetric}.

\section{Validation}

The Catch2 suite (\texttt{test\_kerr\_schild.cpp}) checks the following:
\begin{itemize}
\item \eqref{eq:ksmetric} reproduces the BL Kerr metric, transformed, to
$10^{-12}$ at 240 points;
\item the inverse and the analytic gradient match their numerical
counterparts;
\item the Schwarzschild limit gives the critical impact parameter
$\sqrt{27}\,M$;
\item the constants of motion are conserved;
\item the camera traces the photon it receives (thesis Appendix on image
handedness).
\end{itemize}
Against the thesis' Carter tracer the image-plane comparison of the thesis
appendix gives, over $\HandKSPixels$ pixels off the spin-axis column:
\begin{itemize}
\item disk-hit radii agree to $\HandKSOffAxisDr$ and redshift factors to
$\HandKSOffAxisDg$;
\item the worst drifts along the Kerr--Schild rays are $\HandKSConsL$ in
$L_z$, $\HandKSConsQ$ in $Q$ and $\HandKSConsH$ in $H$, far inside the
thesis' $10^{-6}$ health criterion.
\end{itemize}
On the image's central column the agreement is $\HandKSAxisDr$. That column
is where BL rays pass closest to the polar coordinate singularity, so the
difference measures the reference, not the new integrator.

The time coordinate needs its own check once sources vary (Chapter~\ref{ch:xflares}).
Along a Kerr--Schild ray $dt/d\lambda = -p_t + f\,(l\cdot p)$. In the outgoing
chart the KS time of an equatorial crossing must equal the BL time from
Carter's $t$-equation plus $s\,T(r)$. The flare test suite verifies this to
$10^{-7}\,M$ for three rays of the GRAVITY flare geometry. Building that test
taught a lesson about the reference. A fresh Carter integration without the
thesis' constraint projection, which re-imposes $p_r^2 = R(r)$ and
$p_\theta^2 = \Theta(\theta)$ after every step, agreed with the Kerr--Schild
time only to $10^{-4}$ at this near-face-on inclination. It converged to the
Kerr--Schild value only as its tolerance was tightened, whereas the
Kerr--Schild result was unchanged to twelve digits at every tolerance. With
the projection restored, the two agree to $10^{-10}$.

\section{The live kernel}

The GPU kernel is a single-precision port of the twin:
\begin{itemize}
\item classical fourth-order Runge--Kutta with steps of a few per cent of
$r$, finer inside the photon region and near the disk plane;
\item bisection of each equatorial crossing on the step fraction;
\item capture at $r < r_+$;
\item for escaping rays, the asymptotic direction from the coordinate
velocity, so the background map has no parallax.
\end{itemize}
The seam heuristics of the BL kernel are gone, and the axis column renders
like any other. Three quantitative comparisons with fp64 references
underwrite it:
\begin{itemize}
\item the colour comparison of \S\ref{sec:livecolour}: $\XLiveBeamErr$\,\%
in Doppler contrast;
\item the moving-observer shadow edges of Chapter~\ref{ch:xobs}:
$\XObsGpuMaxPx$\,pixel;
\item the flare images of Chapter~\ref{ch:xflares}: centroids within
$\XFlGpuCentroidPx$\,pixel, total flux within $\XFlGpuFluxDev$\,\%.
\end{itemize}
